\documentclass[11pt,letterpaper]{article}

\usepackage[margin=1in]{geometry}
\usepackage[T1]{fontenc}
\usepackage[utf8]{inputenc}
\usepackage{lmodern}
\usepackage{microtype}
\usepackage{xcolor}
\usepackage{colortbl}
\usepackage{array}
\usepackage{booktabs}
\usepackage{longtable}
\usepackage{tabularx}
\usepackage{enumitem}
\usepackage{xparse}
\usepackage{fancyhdr}
\usepackage{setspace}
\usepackage[hidelinks]{hyperref}

\setlength{\parindent}{0pt}
\setlength{\parskip}{0.65em}
\setlength{\headheight}{14pt}
\setlist[itemize]{leftmargin=1.5em,itemsep=0.25em,topsep=0.25em}
\setlist[enumerate]{leftmargin=1.7em,itemsep=0.25em,topsep=0.25em}
\renewcommand{\arraystretch}{1.18}

\pagestyle{fancy}
\fancyhf{}
\lhead{RFE Response -- EB-1A Petition}
\rhead{Yu Wang, Ph.D.}
\cfoot{\thepage}

% ---------------------------------------------------------------------------
% Editable case metadata
% ---------------------------------------------------------------------------
\newcommand{\LetterDate}{[Month Day, Year]}
\newcommand{\Beneficiary}{Yu Wang, Ph.D.}
\newcommand{\Classification}{INA \S 203(b)(1)(A)}
\newcommand{\PetitionType}{Form I-140}
\newcommand{\Field}{Transportation Engineering}
\newcommand{\IntendedRole}{Lead Autonomy Engineer}
\newcommand{\USCISAddress}{%
USCIS\\
Attn: I-140 (Box 660128)\\
2501 S. State Hwy. 121 Business\\
Suite 400\\
Lewisville, TX 75067-8003%
}

% ---------------------------------------------------------------------------
% Exhibit-reference helpers
%
% Edit only the text inside the first braces:
%   \ExhibitRef{R-1}
%   \ExhibitRef{12B}
%   \ExhibitsRef{4}{R-2}
%   \Exhibit{R-1}{Officer RFE notice}
%   \Exhibits{4}{R-2}{Expert letter and corroborating record}
%
% Single exhibit IDs are detokenized, so suffixes like R-1 and 12B print
% correctly. For two exhibits at the same time, use the plural helpers with
% two IDs instead of typing & inside one ID.
% ---------------------------------------------------------------------------
\NewDocumentCommand{\ExhibitID}{m}{\detokenize{#1}}
\NewDocumentCommand{\ExhibitRef}{m}{\textbf{Exhibit~\ExhibitID{#1}}}
\NewDocumentCommand{\ExhibitsRef}{m m}{\textbf{Exhibits~\ExhibitID{#1}\&\ExhibitID{#2}}}
\NewDocumentCommand{\Exhibit}{m m}{\ExhibitRef{#1}: #2}
\NewDocumentCommand{\Exhibits}{m m m}{\ExhibitsRef{#1}{#2}: #3}
\NewDocumentCommand{\TODO}{m}{\textcolor{red}{[TODO: #1]}}

\newcolumntype{Y}{>{\raggedright\arraybackslash}X}
\newcommand{\ResponseTableHeader}{\rowcolor{black!8}}

\begin{document}
\doublespacing

\LetterDate

\vspace{1em}
\USCISAddress

\vspace{1.5em}
\textbf{RE: Response to Request for Evidence for Permanent Residency Petition for Alien of Extraordinary Ability (EB-1A)}

\begin{tabular}{@{}ll@{}}
\textbf{Petitioner/Beneficiary:} & \Beneficiary \\
\textbf{Classification Sought:} & \Classification \\
\textbf{Type of Petition:} & \PetitionType \\
\textbf{Field:} & \Field \\
\textbf{Intended Role:} & \IntendedRole \\
\end{tabular}

\vspace{1em}
Dear USCIS Officer:

I respectfully submit this response to the Request for Evidence (RFE) issued in connection with my \PetitionType{} Immigrant Petition for Alien of Extraordinary Ability (EB-1A). I appreciate the opportunity to address the officer's concerns directly and to provide additional objective evidence in support of my petition.

At the outset, I note and appreciate that USCIS has already accepted two regulatory criteria:

\begin{itemize}
  \item \textbf{Judging the work of others:} USCIS found that sufficient evidence was submitted to show that I have participated as a judge of the work of others in the field of \Field.
  \item \textbf{Authorship of scholarly articles:} USCIS found that sufficient evidence was submitted to show authorship of scholarly articles in the field, in professional or major trade publications or other major media.
\end{itemize}

This response therefore focuses on the criterion that the officer has specifically requested additional evidence to evaluate: \textbf{original contributions of major significance in the field}. This response also addresses \textbf{final merits}, showing that the totality of my record---including the two accepted criteria, the original-contributions evidence submitted below, and my sustained recognition across academia, industry, and government---demonstrates that I am among the small percentage at the very top of the field.

In this response, I do not re-argue authorship or judging. Instead, I supply \textbf{objective documentary evidence}, not merely support letters, to establish the following:

\begin{enumerate}
  \item \textbf{My specific personal contributions were original.} I identify two field-level contributions and explain, with corroborating evidence, what I personally created that did not exist before.
  \item \textbf{My contributions have been used, adopted, extended, benchmarked, or relied upon} by independent researchers, government-facing transportation professionals, academic collaborators, and industry practitioners.
  \item \textbf{My contributions are of major significance} because they have moved beyond publication and into practical use across the communities that matter most in autonomous driving: academia, industry, and government.
\end{enumerate}

The officer's RFE helpfully identifies the types of evidence that would be probative---including adoption, implementation, or use by others. I have organized this response to address those concerns directly. The evidence is presented in two technical contribution sections, each following a consistent structure that identifies:

\begin{itemize}
  \item The problem that existed before my work;
  \item What I created and why it was original;
  \item My personal role, supported by objective attribution evidence;
  \item Objective evidence of adoption, use, or reliance by independent users;
  \item Cross-sector translation evidence showing practical impact; and
  \item Why this contribution demonstrates major significance in the field.
\end{itemize}

Following those two contribution sections, I present an agency-by-agency rebuttal table directly addressing each entity the officer specifically identified, supplying the objective documentary evidence requested. I then conclude with a final merits argument and a complete exhibit cross-reference list.

I do not present panel discussions, NSF support, Wisconsin DOT-facing training, advisory roles, or patents as separate contributions. Instead, I use them as objective evidence of translation, adoption, reliance, and practical use supporting my two technical original contributions.

I respectfully submit that the evidence provided in this response, together with the record already accepted by USCIS, establishes original contributions of major significance and sustained acclaim at the top of the field. I respectfully request approval of my petition.

\tableofcontents
\newpage

\section{Introduction and Summary of the RFE Response}

\subsection{Overview of the RFE Context}

USCIS has accepted that I satisfy the judging criterion and the authorship criterion. The central issue in this response is whether my work satisfies the original-contributions criterion and, in the final merits analysis, whether the totality of my record shows sustained acclaim and recognition at the top of the field.

The evidence below is organized around two technical original contributions:

\begin{enumerate}
  \item \textbf{Stochastic and empirical modeling of human-driver behavior in mixed AV-human traffic.}
  \item \textbf{Optimization- and learning-based trajectory planning/control methods for mixed traffic under human-driver disruption and real-world deployment constraints.}
\end{enumerate}

These contributions form a coherent technical arc. First, I advanced how the field understands and models uncertain human-driver behavior in mixed AV-human traffic. Second, I advanced how AVs and CAVs can plan and control their movement in that mixed traffic.

These contributions are not presented as mere publications. My publications disclosed the contributions, but the contributions themselves are the underlying models, datasets, empirical findings, analytical frameworks, control methods, benchmarks, and autonomous-driving technologies that others later adopted, extended, benchmarked, or relied upon.

\subsection{Structure of This Response}

This response is structured to address the officer's concerns directly and provide the objective documentary evidence requested:

\begin{enumerate}
  \item \textbf{Two technical contribution sections} (Contributions 1--2), each following a consistent structure that identifies: the problem that existed before my work; what I created and why it was original; my personal role, supported by objective attribution evidence; objective evidence of adoption, use, or reliance by independent users; cross-sector translation evidence integrated within each contribution; and why this contribution demonstrates major significance in the field.
  \item \textbf{Agency-by-agency rebuttal table} directly addressing each entity the officer specifically identified, supplying the objective documentary evidence requested.
  \item \textbf{Final merits argument} demonstrating that the totality of my record establishes sustained acclaim at the top of the field.
  \item \textbf{Exhibit cross-reference list} with all exhibits organized by contribution and evidence type.
\end{enumerate}

I respectfully submit that the evidence provided in this response, together with the record already accepted by USCIS, establishes original contributions of major significance and sustained acclaim at the top of the field. I respectfully request approval of my petition.

\section{Roadmap of the Original Contributions}

\begin{tabularx}{\textwidth}{@{}p{0.20\textwidth}YYY@{}}
\toprule
\textbf{Contribution} & \textbf{Specific Original Contribution} & \textbf{Main Technical Evidence} & \textbf{Translation / Reliance Evidence} \\
\midrule
Contribution 1 & Stochastic and empirical modeling of human-driver behavior in mixed AV-human traffic. & First-authored stochastic car-following model; controlled HV-following-AV field experiment; mixed-traffic dataset; AV appearance / trust finding. & NSF / academic project support; Wisconsin DOT-facing training; AMS/CUTR/FHWA-related use; regulatory/policy discussion; independent researchers extending the work. \\
\addlinespace
Contribution 2 & Optimization- and learning-based trajectory planning/control for mixed traffic under human-driver disruption and deployment constraints. & Trajectory optimization papers; proactive longitudinal control papers; reinforcement-learning methods; patents on AV learning, vehicle action prediction, and intention prediction. & FHWA/CARMA/Leidos-type use; multi-stakeholder panel; industry adoption; benchmark use; patents as practical translation evidence. \\
\bottomrule
\end{tabularx}

\vspace{1em}
These two contributions form a coherent technical arc. First, I advanced how the field understands and models uncertain human-driver behavior in mixed AV-human traffic. Second, I advanced how AVs and CAVs can plan and control their movement in that mixed traffic. The cross-sector evidence does not stand alone as a separate contribution; instead, it shows that my technical contributions were valued, relied upon, and translated beyond ordinary publication.

\section{Clarification: Publications Are Not the Contributions Themselves}

USCIS has already accepted the authorship criterion. I understand that the present RFE concerns a separate requirement: whether my work constitutes original contributions of major significance in the field. I therefore do not rely on publication alone to establish this criterion.

I respectfully do not claim that my publications alone establish original contributions of major significance. My peer-reviewed articles are the vehicles through which my work was disclosed to the field. The claimed contributions are the underlying models, datasets, empirical findings, analytical frameworks, control methods, benchmarks, and autonomous-driving technologies that independent researchers, agencies, industry practitioners, and academic collaborators later used, adopted, extended, benchmarked, or relied upon.

The proof of major significance comes from the independent use, practical reliance, and cross-sector influence of the underlying work---not from publication count, citation count, or support letters standing alone.

\section{Contribution 1: Stochastic and Empirical Modeling of Human-Driver Behavior in Mixed AV-Human Traffic}

\textbf{Technical subtitle:} Generalized stochastic car-following framework, controlled HV-following-AV field experiment, mixed-traffic behavioral dataset, and AV appearance / behavioral adaptation finding.

\subsection{Problem Before My Work}

Autonomous vehicles will share roads with human drivers for many years. Researchers, engineers, and government agencies therefore need realistic ways to model how human drivers behave in mixed AV-human traffic.

Before my work, the field faced two related gaps:

\begin{enumerate}
  \item Many car-following models treated human driving behavior as deterministic or did not adequately capture uncertainty, variation, and heterogeneity in human behavior.
  \item Researchers lacked controlled real-world evidence showing whether human drivers behave differently when following an autonomous vehicle rather than a human-driven vehicle.
\end{enumerate}

A few key terms: \textbf{mixed traffic} means roads where autonomous vehicles and human-driven vehicles operate together. A \textbf{car-following model} predicts how a vehicle follows the vehicle ahead. \textbf{Stochastic} means accounting for uncertainty and variation. A \textbf{controlled field experiment} is a real-world experiment designed to observe a specific behavior under controlled conditions, rather than passively collecting ordinary traffic data.

\subsection{What I Created and Why It Was Original}

I contributed a more realistic modeling foundation for mixed AV-human traffic by combining:

\begin{enumerate}
  \item A generalized stochastic car-following framework that accounts for uncertainty and variation in human driving behavior; and
  \item Controlled empirical evidence showing how human drivers behave when following an autonomous vehicle, including the finding that AV recognizability and subjective trust affect driver behavior.
\end{enumerate}

\TODO{Add specific technical description of the stochastic model formulation, the field experiment design, and the AV appearance / trust finding.}

The originality is not that I worked on ``traffic simulation'' generally. The originality is that I addressed a specific missing link in mixed-traffic modeling: how to represent uncertain human-driver behavior and how to incorporate observed behavioral adaptation to AVs.

\subsection{My Personal Role}

\begin{tabularx}{\textwidth}{@{}p{0.28\textwidth}p{0.36\textwidth}Y@{}}
\toprule
\textbf{Work / Evidence} & \textbf{My Specific Role} & \textbf{Exhibit} \\
\midrule
First-authored stochastic car-following paper & \TODO{Formulated model, developed analytical framework, conducted analysis, wrote manuscript.} & \ExhibitRef{[ID]} \\
Field experiment on HV-following-AV behavior & \TODO{Experiment design, data processing, analysis, interpretation, dataset dissemination.} & \ExhibitRef{[ID]} \\
Coauthor / corresponding-author statement & \TODO{Confirms specific personal contribution rather than team attribution.} & \ExhibitRef{[ID]} \\
Dataset / experiment documentation / dissemination & \TODO{Shows contribution was not merely a publication but an underlying empirical resource.} & \ExhibitRef{[ID]} \\
\bottomrule
\end{tabularx}

\subsection{Objective Evidence of Adoption, Use, Extension, or Reliance}

\begin{tabularx}{\textwidth}{@{}p{0.18\textwidth}p{0.18\textwidth}p{0.28\textwidth}Y@{}}
\toprule
\textbf{Independent User / Entity} & \textbf{What They Used} & \textbf{How They Used It} & \textbf{Exhibit} \\
\midrule
Independent researchers extending HV-following-AV experiments & Field experiment protocol, dataset, or AV appearance finding & Extended to T-intersections or different recognizability settings. & \ExhibitRef{[ID]} \\
CUTR / FHWA-related AMS work & Mixed-traffic data or behavioral findings & Used to calibrate or modify mixed-traffic simulation tools. & \ExhibitRef{[ID]} \\
Texas A\&M / later mixed-traffic researchers & Stochastic model format or modeling approach & Used as foundation for later CAV longitudinal control methods. & \ExhibitRef{[ID]} \\
Virginia / VTRC & AV behavior evidence or car-following findings & Used in safety, automation, or AV-human interaction analysis. & \ExhibitRef{[ID]} \\
European Commission / JRC & Human-driver heterogeneity or AV behavior findings & Used in driver heterogeneity, policy, or regulatory analysis. & \ExhibitRef{[ID]} \\
Policy / regulatory researchers & AV appearance or recognizability finding & Used in discussions of behavioral adaptation or AV appearance regulation. & \ExhibitRef{[ID]} \\
\bottomrule
\end{tabularx}

\subsection{Cross-Sector Translation Evidence}

The following evidence is not presented as a separate contribution. It shows that government-facing and academic stakeholders sought my technical expertise beyond publication.

\subsubsection{NSF / Academic Project Support}

\TODO{Include if evidence shows academic researchers sought my industry-side input.}

\TODO{Insert framing: academic researchers did not merely cite my published work; they sought my industry-grounded expertise to connect mixed-traffic research with real-world autonomous-driving deployment constraints, data limitations, and validation needs.}

\TODO{Insert evidence: NSF proposal excerpt, PI letter, email requesting support, acknowledgment or project document, evidence tying the project to mixed traffic, AV-human interaction, simulation, validation, deployment, or data.}

\subsubsection{Wisconsin DOT-Facing Training}

\TODO{Include here if the training topic relates to AV behavior, simulation, modeling, safety, data, or deployment constraints.}

\TODO{Insert framing: government-facing transportation professionals sought my technical expertise in mixed-traffic behavior, AV modeling, simulation, safety, and deployment constraints.}

\TODO{Insert evidence: training agenda, speaker list, organizer letter, audience description, slides or topic outline, proof that I was the only or key industry representative.}

\subsection{Why Contribution 1 Shows Major Significance}

\TODO{Conclude: independent researchers and government-related projects did not merely cite my work as background. They used my stochastic modeling framework, empirical data, field-experiment protocol, behavioral findings, and mixed-traffic insights to improve simulation, extend experiments, guide later CAV control research, and analyze human-driver adaptation around autonomous vehicles. This evidence shows substantive reliance on my contribution rather than ordinary publication or citation.}

\subsection{Supporting Expert Letters}

\TODO{Use expert letters only after objective evidence. Reference: Yang Zhou (stochastic model originality, benchmark value, field experiment significance); Jianming Ma (policy/regulatory significance, AV appearance, mixed-traffic data, DOT relevance); coauthor/corresponding author (personal attribution); NSF/project organizer (reliance on industry-grounded expertise, if available); Wisconsin DOT training organizer (government-facing reliance, if available).}

\section{Contribution 2: Optimization- and Learning-Based Trajectory Planning/Control for Mixed Traffic Under Human-Driver Disruption and Deployment Constraints}

\textbf{Technical subtitle:} Trajectory optimization, proactive longitudinal control, reinforcement-learning trajectory control, transfer learning, and patented AV prediction / planning / learning methods.

\subsection{Problem Before My Work}

Autonomous and connected vehicles must operate around human-driven vehicles that may cut in, change lanes, slow down unpredictably, or disrupt traffic flow. Earlier AV/CAV planning and control methods often relied on simplified assumptions, fixed input structures, or idealized traffic behavior. These assumptions limited their usefulness for real-world mixed traffic, where AVs must make safe, efficient, smooth, and reliable decisions around human drivers.

A few key terms: \textbf{trajectory optimization} means choosing a vehicle's speed and path to improve safety, travel time, fuel use, or comfort. \textbf{Reinforcement learning} is an AI method where a system learns decisions by testing actions and receiving rewards or penalties. \textbf{Proactive longitudinal control} means controlling vehicle motion in advance to manage or prevent disruptions caused by human-driven vehicles. A \textbf{benchmark} is a reference method or dataset that others use to test whether their own systems work well. A \textbf{deployment constraint} is a real-world limitation such as missing data, variable input size, safety requirements, computational limits, or unpredictable human behavior.

\subsection{What I Created and Why It Was Original}

I developed optimization- and learning-based methods for trajectory planning and control in mixed traffic, especially addressing the practical problem that human-driven vehicles disrupt AV/CAV operations through uncertain behavior and lane changes.

My work advanced the field through:

\begin{enumerate}
  \item Trajectory optimization methods for safety, efficiency, smoothness, and energy/fuel performance.
  \item Proactive longitudinal control methods to manage or preclude disruptive lane changes by human-driven vehicles.
  \item Reinforcement-learning approaches that improve trajectory control under realistic data and deployment constraints.
  \item Patented AV prediction, planning, and learning methods that show continued translation of this research line into implementable AV technology.
\end{enumerate}

\TODO{Add specific technical description of the optimization formulation, proactive control logic, RL approach with transfer learning, transfer parameter selection, data imputation, variable-input-size handling, etc.}

\subsection{My Personal Role}

\begin{tabularx}{\textwidth}{@{}p{0.30\textwidth}p{0.34\textwidth}Y@{}}
\toprule
\textbf{Work / Evidence} & \textbf{My Specific Role} & \textbf{Exhibit} \\
\midrule
Trajectory optimization papers & \TODO{Model design, optimization framework, benchmarking, simulation, analysis.} & \ExhibitRef{[ID]} \\
Proactive longitudinal control / disruptive lane-change papers & \TODO{Method development, control logic, experiment design, evaluation.} & \ExhibitRef{[ID]} \\
Reinforcement-learning papers & \TODO{RL formulation, data imputation, transfer parameter selection, evaluation.} & \ExhibitRef{[ID]} \\
Granted U.S. patents (RL transfer learning; hierarchical vehicle action prediction; intention prediction in symmetric scenarios) & Named inventor on AV learning, action prediction, and intention prediction inventions. & \ExhibitRef{[ID]} \\
Employer / manager / collaborator confirmation & \TODO{Confirm personal role in AV prediction, planning, control, or implementation.} & \ExhibitRef{[ID]} \\
\bottomrule
\end{tabularx}

\subsection{Objective Evidence of Adoption, Use, Extension, or Reliance}

\begin{tabularx}{\textwidth}{@{}p{0.18\textwidth}p{0.22\textwidth}p{0.28\textwidth}Y@{}}
\toprule
\textbf{Independent User / Entity} & \textbf{What They Used} & \textbf{How They Used It} & \textbf{Exhibit} \\
\midrule
FHWA / CARMA-related evidence & Trajectory optimization methods or benchmarks & Used as foundation for operational concepts, simulation evaluation, or CAV research. & \ExhibitRef{[ID]} \\
Leidos / FHWA-related engineering & Trajectory benchmarks, field data, or simulation checks & Used to evaluate simulation reliability, smoothness, safe gaps, fuel use, or safety behavior. & \ExhibitRef{[ID]} \\
Independent researchers & Proactive control, RL, or trajectory optimization methods & Extended or benchmarked against my methods. & \ExhibitRef{[ID]} \\
Pony.ai / other AV company & Mixed-traffic control framework or lane-change assistance concept & Used to improve mixed-traffic lane-change algorithms or safety evaluation. & \ExhibitRef{[ID]} \\
Book / editorial invitation & Expertise in AI, RL, or trajectory optimization & Invited to edit or contribute to a major technical book or chapter. & \ExhibitRef{[ID]} \\
\bottomrule
\end{tabularx}

\subsection{Patents as Practical Translation Evidence}

I do not rely on patents alone to establish the original-contributions criterion. Rather, my granted U.S. patents corroborate that my research expertise has continued to translate into practical autonomous-vehicle prediction, planning, and learning systems. These patents support the broader record showing that my contributions are not confined to academic publications but extend into implementable AV technologies.

\begin{tabularx}{\textwidth}{@{}p{0.28\textwidth}p{0.32\textwidth}Y@{}}
\toprule
\textbf{Patent} & \textbf{Technical Relevance} & \textbf{Exhibit} \\
\midrule
Reinforcement learning technique for network-based transfer learning & RL transfer learning / target-domain adaptation---supports learning-based AV planning/control expertise under deployment constraints. & \ExhibitRef{[ID]} \\
Hierarchical vehicle action prediction & Vehicle action prediction / planning stack---supports AV prediction and planning for complex traffic scenarios. & \ExhibitRef{[ID]} \\
Intention prediction in symmetric scenarios & Intention prediction / scenario understanding---supports practical AV behavior prediction and planning. & \ExhibitRef{[ID]} \\
\bottomrule
\end{tabularx}

\subsection{Cross-Sector Translation Evidence}

The following evidence is not presented as a separate contribution. It shows that government, industry, and multi-stakeholder groups sought my technical expertise in AV planning, prediction, safety validation, and deployment beyond ordinary academic publication.

\subsubsection{Multi-Stakeholder Panel}

\TODO{Include if the panel topic relates to planning, prediction, deployment, AV safety validation, testing, industry constraints, or commercialization.}

\TODO{Insert framing: stakeholders from academia, government, and industry sought my real-world AV expertise in planning, prediction, safety validation, deployment constraints, and commercialization.}

\TODO{Insert evidence: panel agenda or program, speaker list and affiliations, organizer letter, panel description, proof that I was the only AV industry practitioner if accurate, slides or notes, audience description.}

\subsubsection{Industry / Government-Facing Implementation}

\TODO{Use evidence from Leidos, FHWA, CARMA, AMS, or similar projects here.}

\TODO{Insert framing: my trajectory optimization and learning-based control work was used or relied upon in government-related and industry-facing AV/CAV simulation, operational-concept, or benchmarking contexts.}

\TODO{Insert evidence: FHWA or CARMA report excerpts, project documents, public technical reports, citation-context excerpts, Leidos/project manager letter tied to objective documents.}

\subsection{Why Contribution 2 Shows Major Significance}

\TODO{Conclude: my trajectory optimization, proactive control, and learning-based methods were used as benchmarks, extended by independent researchers, and connected to government-related and industry-facing AV/CAV simulation and deployment contexts. My patents further corroborate that the same research expertise translated into implementable AV prediction, planning, and learning technologies. This evidence shows practical reliance on my technical work, not merely recognition of my publications.}

\subsection{Supporting Expert Letters}

\TODO{Use expert letters only after objective evidence. Reference: Pawel Gora (reinforcement learning, benchmark value, commercialization relevance, missing-data/variable-input-size issue); Amir Ghiasi (FHWA/CARMA/Leidos implementation and simulation benchmark value, if supported by objective documents); Jianming Ma (traffic operations, trajectory optimization, speed management, CARMA/FHWA relevance); employer/manager/collaborator (personal role in AV prediction, planning, control, or implementation); panel organizer (why my industry expertise was sought, if available).}

\section{Agency-by-Agency and Entity-by-Entity Rebuttal}

This table directly addresses the officer's concern that prior claims about government or agency use were not objectively documented. It cross-references the evidence presented under each contribution and shows exactly where each challenged claim is supported.

\begin{longtable}{@{}p{0.15\textwidth}p{0.18\textwidth}p{0.22\textwidth}p{0.18\textwidth}p{0.17\textwidth}@{}}
\toprule
\textbf{Agency / Entity} & \textbf{Officer Concern} & \textbf{Objective Evidence Submitted} & \textbf{What Was Used} & \textbf{Related Contribution} \\
\midrule
\endhead
U.S. DOT / FHWA & Prior claim lacked objective proof & \TODO{FHWA report, CARMA report, AMS report, or citation-context excerpt.} & \TODO{Trajectory optimization method, benchmark, field data, mixed-traffic findings.} & Contribution 1 or 2 \\
Florida / CUTR & Claimed use requires corroboration & \TODO{AMS / CUTR report or paper excerpt.} & \TODO{Mixed-traffic data or findings.} & Contribution 1 \\
Texas / TxDOT & Policy impact needs objective support & \TODO{TxDOT-related document or technical evidence, plus expert letter.} & \TODO{Traffic operations, speed management, mixed-traffic work.} & Contribution 1 or 2 \\
Virginia / VTRC & Claimed use needs support & \TODO{VTRC report or citation-context excerpt.} & \TODO{AV behavior evidence, car-following model, safety analysis.} & Contribution 1 \\
European Commission / JRC & International relevance needs proof & \TODO{EC / JRC report or paper excerpt.} & \TODO{Human-driver heterogeneity or AV behavior findings.} & Contribution 1 \\
Wisconsin DOT-facing training & Government-facing use should be documented & \TODO{Agenda, organizer letter, slides, speaker list.} & \TODO{AV modeling, simulation, mixed-traffic behavior, deployment expertise.} & Contribution 1 / final merits \\
Leidos / FHWA-related projects & Industry/government implementation needs objective proof & \TODO{Project report, public document, CV/project role evidence, or letter tied to report.} & \TODO{Trajectory benchmark, field data, simulation method.} & Contribution 2 \\
Pony.ai / other companies & Company use needs proof & \TODO{Company paper, technical document, citation-context excerpt.} & \TODO{Mixed-traffic control framework, field experiment protocol, dataset.} & Contribution 1 or 2 \\
\bottomrule
\end{longtable}

\section{Final Merits}

The totality of the evidence demonstrates sustained national and international acclaim and shows that I am among the small percentage who have risen to the very top of the field.

\subsection{USCIS Has Already Accepted Two Independent Criteria}

USCIS has accepted that I satisfy the authorship criterion and the judging criterion. These two criteria, standing alone, already provide independent evidence of my standing in the field.

\subsection{The Original-Contributions Evidence Shows Field-Level Impact}

The evidence submitted above establishes a third independent criterion---original contributions of major significance---by showing that my two technical contributions were not merely published but were used, adopted, extended, benchmarked, and relied upon by independent researchers, government-facing professionals, and industry practitioners:

\begin{enumerate}
  \item \textbf{Contribution 1:} I advanced how the field models human-driver behavior in mixed AV-human traffic, through a generalized stochastic car-following framework, a controlled field experiment, a mixed-traffic behavioral dataset, and a finding about AV appearance and behavioral adaptation. Independent researchers extended this work, government-related simulation projects used it for calibration, and policy researchers used it in regulatory analysis.
  \item \textbf{Contribution 2:} I advanced how AVs and CAVs can plan and control movement under mixed-traffic disruptions and deployment constraints, through trajectory optimization, proactive longitudinal control, reinforcement-learning methods, and patented AV prediction/planning/learning technologies. Government-related and industry-facing projects relied on these methods as benchmarks, and multi-stakeholder panels sought my expertise for deployment and safety discussions.
\end{enumerate}

\subsection{My Recognition Is Sustained and Cross-Sector}

My record shows sustained recognition across academia, industry, and government:

\begin{itemize}
  \item Independent academic use of my models, data, and methods.
  \item Government-facing use in simulation, training, and policy analysis.
  \item Industry-facing use through benchmarks, patents, and engineering projects.
  \item Multi-stakeholder panel participation as the only AV industry practitioner.
  \item NSF / academic project support showing reliance on my industry-grounded expertise.
  \item Wisconsin DOT-facing training showing government-seeking of my technical knowledge.
  \item Advisory and reviewing roles.
  \item Lead Autonomy Engineer position.
\end{itemize}

\subsection{National Benefit and Continued Work}

I will continue working as a Lead Autonomy Engineer and transportation engineering researcher, contributing to AV safety, AV reliability, traffic efficiency, mixed-traffic modeling, deployment readiness, and U.S. technological leadership in autonomous vehicles.

\subsection{Conclusion}

Considering the accepted authorship and judging criteria together with the objective evidence of my original technical contributions, their independent adoption and use, their practical translation into government-facing and industry-facing contexts, and my sustained recognition across academia, industry, and government, the record establishes that I satisfy the original-contributions criterion and meet the final merits standard for EB-1A classification.

\section{Exhibit Cross-Reference List}

\begin{longtable}{@{}p{0.22\textwidth}p{0.70\textwidth}@{}}
\toprule
\textbf{Exhibit ID} & \textbf{Description} \\
\midrule
\endhead
\ExhibitRef{R-1} & \TODO{RFE notice or response cover evidence.} \\
\ExhibitRef{12B} & \TODO{Example editable exhibit ID with letter suffix.} \\
\ExhibitsRef{4}{R-2} & \TODO{Example two-exhibit reference.} \\
\ExhibitRef{[ID]} & \TODO{Add each exhibit used in the response.} \\
\bottomrule
\end{longtable}

\vspace{2em}
Respectfully submitted,

\vspace{3em}
\Beneficiary

\end{document}
